\chapter{Compilação e Geração do PDF}

\section{Fluxo automatizado}

Na raiz do projeto, execute:

\begin{lstlisting}
latexmk -outdir=build/perfis
\end{lstlisting}

O arquivo \texttt{.latexmkrc} seleciona LuaLaTeX, aciona o Biber quando necessário e define \texttt{tcc.tex} como entrada padrão. Para o exemplo completo:

\begin{lstlisting}
latexmk -outdir=build/perfis tcc-exemplo-undf.tex
\end{lstlisting}

\section{Fluxo manual}

Quando o \texttt{latexmk} não estiver disponível, use:

\begin{lstlisting}
New-Item -ItemType Directory -Force build/perfis | Out-Null
lualatex -output-directory=build/perfis tcc.tex
biber --input-directory=build/perfis --output-directory=build/perfis tcc
lualatex -output-directory=build/perfis tcc.tex
lualatex -output-directory=build/perfis tcc.tex
\end{lstlisting}

No PowerShell, as saídas ficam em \texttt{build/perfis/}. O script
\texttt{scripts/verificar-perfis.ps1} automatiza esse ciclo no Windows, inclusive
sem Perl. O Biber deve receber o nome-base, sem a extensão \texttt{.tex}.
Passagens adicionais podem ser necessárias após alterações em sumário, listas,
citações ou referências cruzadas.

\section{Overleaf}

Envie a estrutura completa do projeto e selecione LuaLaTeX como compilador. O arquivo \texttt{.latexmkrc} acompanha o repositório. Confirme no painel de logs se a bibliografia foi processada e se não restaram referências indefinidas.

\section{Arquivos auxiliares}

Arquivos como \texttt{.aux}, \texttt{.bcf}, \texttt{.bbl}, \texttt{.toc}, \texttt{.lof}, \texttt{.lot} e \texttt{.lofl} registram dados entre as passagens. Eles não fazem parte do texto autoral e podem ser regenerados. Para limpar artefatos pelo \texttt{latexmk}, use \texttt{latexmk -outdir=build/perfis -c}; faça nova compilação completa em seguida.

\section{Critério de conclusão}

Considere a geração tecnicamente estabilizada somente quando: o comando terminar sem erro fatal; o PDF abrir normalmente; não houver pedido de nova execução do Biber; sumário, listas e referências cruzadas estiverem resolvidos; e a inspeção visual das páginas críticas tiver sido feita.
